\scanpage{48}
\margin{Avril 86}
\uline{Trialités}\qquad Soit $C$ catégorie, $n$ entier naturel $\geq\deleted{\unsure{3}}2$.
\inserted{(le cas $n=2$ trivial)}
On appelle $n$-alité (trialité, quadralité, pentalité si $n=3,4,5$)
un système
\[
\{I,E,\{E_i\}_{i\in I},(\varphi_i)_{(i\in I)}\}
\]
où \quad $I$ \deleted{\unsure{ens|ens.}} de cardinal $n$,
\[
\varphi_i:E\mathrel{\unsure{-|\to}}E_i\qquad i\in I
\]
flèches de $C$, satisfaisant

$\forall J\subset I$ avec $\operatorname{card}J=2$, l'application
\[
E\xrightarrow{\varphi_J}\underbrace{E_J=\prod_{i\in J}E_i}_{
\substack{\deleted{\text{on suppose que}}\\\deleted{\text{ce produit existe,}}\\
\deleted{\forall J\in\unsure{\mathfrak P_2|\mathcal P_2}(I)}}}
\tag*{(1)}
\]
est un \uline{iso} (dans $\widehat C$) $E$ entouré des $\varphi_i$ $(i\in J)$ est un produit des
$E_i$ dans $C$.

\margin{[\, Les $n$-alités forment une catégorie $n$-al$(C)$, fibrée sur le groupoïde \deleted{\unsure{Ens}}\unsure{$\mathcal S_n$|$\mathrm{Ens}_n$}}

Je m'intéresse surtout \deleted{aux} deux cas

$1^\circ)$ $C$ catégorie additive

$2^\circ)$ $C=\mathrm{Ens}$

\deleted{Dans le cas \unsure{2}} Je vais ici, regarder le cas $2^\circ)$

Si un des $E_i$ est vide, alors \inserted{$E$ l'est et} tous sauf au plus
un le sont. On va exclure ce cas, \inserted{dit} ``impropre''.

Dans tous les autres cas les $E_i$ peuvent
être considérés comme des quotients de $E$.
Ainsi, la structure peut être considérée comme
portée par $E$\deleted{\unsure{,}} \inserted{$I$}, \deleted{muni d'une}, une
application $I\mathrel{-}\operatorname{Quot}(E)$. Cette application
est d'ailleurs injective, sauf \deleted{quand \unsure{tous les}}\inserted{dans}
% archive sheet number: 47
